\chapter{Checkpoint Answers}\label{app:answers}

Answers to the \textbf{Checkpoint} questions in each chapter. The Review
Questions at the end of each chapter are deliberately not answered here --- they
are for assessment, and several have more than one defensible answer.

\begin{alertbox}[Use these to check, not to learn]
A checkpoint question you read the answer to teaches almost nothing. Write your
answer down first, then compare. Where yours differs, the useful question is not
``was I right'' but ``what did I not consider''.
\end{alertbox}

\section*{Part I --- Foundations}
\addcontentsline{toc}{section}{Part I --- Foundations}

\subsection*{Chapter 1 --- Networks, Models and This Exam}
\begin{tightnum}
  \item Layers 1 to 3 are ruled out: something reached the server and came back.
        Suspect layer 7 first --- specifically name resolution or the web service
        itself. Test the name (\chref{dns}) and the port (\chref{transport}).
  \item Layer 1. Rising collisions on a link that reports full duplex means a
        duplex mismatch: one end is half duplex and is correctly reporting
        collisions the other end does not see.
  \item Because they are different jobs. A switch's flooding is a \emph{recovery}
        mechanism --- the destination is probably on one of those ports and will
        answer, teaching the switch where it is. A router's unknown network has no
        such prospect; flooding a packet to every interface would loop and
        multiply. So a router drops and reports.
\end{tightnum}

\subsection*{Chapter 2 --- The Physical Layer and Cable Faults}
\begin{tightnum}
  \item That the two ends disagree about duplex. It is certain because a full
        duplex link \emph{cannot} have collisions --- there is no shared medium to
        collide on. The counter can only be incrementing because the local end is
        actually running half duplex, or because the far end is and is causing
        them.
  \item In order: a dirty or damaged connector; the wrong optic or wavelength for
        the distance; and excessive length or a bad splice. Check the cheapest
        first --- clean the connector before ordering hardware.
  \item Because it is unambiguous. Administratively down is a decision somebody
        made, recorded in the configuration, and reversible with one command. A
        port that is up with rising CRC errors is intermittently working, so
        traffic is partly passing, users report ``slow'' rather than ``down'',
        and the cause could be cable, connector, optic, EMI or duplex.
\end{tightnum}

\subsection*{Chapter 3 --- Ethernet, Frames and Switching}
\begin{tightnum}
  \item It drops it. The behaviour is called \term{filtering}: the destination is
        known to be reachable through the same port the frame arrived on, so
        forwarding it would send it back where it came from.
  \item Because the first one triggers an ARP request and possibly a MAC flood,
        and the round trip includes that resolution. Later pings use the cached
        entries.
  \item The MAC address table, on the switch the laptop is connected to. You need
        the laptop's MAC address first --- and note that on wireless it may be
        randomised (\chref{clienttshoot}), so the MAC on the label may not be the
        one on the network.
\end{tightnum}

\subsection*{Chapter 4 --- IPv4 Addressing and Subnetting}
\begin{tightnum}
  \item \cmd{/27} gives a block size of 32 in the fourth octet. The address ends
        in \textbf{37}, which falls in the block starting at 32. Network
        \cmd{10.19.200.32}, broadcast \cmd{10.19.200.63}, first usable
        \cmd{10.19.200.33}, last usable \cmd{10.19.200.62}.
  \item \textbf{No.} \cmd{/25} has a block size of 128, so the host is in
        \cmd{192.168.1.128/25}, whose range is \cmd{.129} to \cmd{.254}. The
        gateway \cmd{192.168.1.126} is in the \emph{other} half,
        \cmd{192.168.1.0/25}. The host will ARP for a gateway that is not on its
        segment and reach nothing beyond its own subnet.
  \item Two. A \cmd{/31} is used on a router-to-router link because there are
        only ever two devices and neither needs a broadcast address, so it
        provides both addresses usefully and halves the address consumption of
        every point-to-point link on a network.
\end{tightnum}

\subsection*{Chapter 5 --- IPv6 Addressing}
\begin{tightnum}
  \item \cmd{2001:db8:0:12::ab}. Leading zeros drop in every group; the single run
        of three all-zero groups becomes \cmd{::}. Note the fourth group is
        \cmd{0}, not part of the \cmd{::} --- only one run may be compressed, and
        the longer one wins.
  \item A global unicast address, which means no router advertisement is being
        received. Look at the router or layer 3 switch for the segment: is
        \cmd{ipv6 unicast-routing} enabled, and does the interface have a global
        address to advertise a prefix from?
  \item Split the MAC, insert \cmd{fffe}, and flip the seventh bit: \cmd{aa}
        becomes \cmd{a8}. The result is \cmd{a8bb:ccff:fe00:0100}.
\end{tightnum}

\subsection*{Chapter 6 --- TCP, UDP and the Diagnostic Toolkit}
\begin{tightnum}
  \item \cmd{U} is an ICMP destination-unreachable reply. Something on the path
        actively told you it could not deliver --- so the packet was routed, and a
        device took responsibility. Five dots means silence, which could be
        anything from a dropped packet to an ACL to a dead host. A reply is
        always more informative than a timeout.
  \item Because a plain ping is sourced from the exit interface's address, which
        the far end may have no route back to. Sourcing from a LAN interface or
        loopback uses an address the far end's routing table knows about.
  \item Any of: an ACL blocking the probes; a device configured not to send ICMP
        time-exceeded messages; asymmetric return path with no route back; a
        firewall dropping the traffic; or the destination genuinely unreachable
        beyond that hop.
\end{tightnum}

\subsection*{Chapter 7 --- The IOS Command Line}
\begin{tightnum}
  \item \cmd{do show ip route}. The \cmd{do} prefix runs a privileged EXEC command
        from any configuration mode without leaving it.
  \item It boots from \cmd{startup-config} in NVRAM. You are editing
        \cmd{running-config} in RAM. Anything not saved is lost at the next
        reload --- see \chref{ansible} for how this catches out automation too.
  \item Because the RSA key pair is named after the fully qualified domain name,
        \cmd{hostname.domain}. Without both, the device has no name to label the
        key with and refuses to generate it.
\end{tightnum}

\section*{Part II --- Switching and Network Access}
\addcontentsline{toc}{section}{Part II --- Switching and Network Access}

\subsection*{Chapter 8 --- VLANs and Access Ports}
\begin{tightnum}
  \item No. VLANs are separate broadcast domains, so the ARP request never reaches
        the other host. Sharing an IP subnet is irrelevant at layer 2 --- and
        because both believe the destination is local, neither will send the
        traffic to a router either. This configuration cannot be fixed by
        routing; the VLANs must be corrected.
  \item \cmd{show interfaces <int> switchport}, which reports both
        \cmd{Administrative Mode} (what you configured) and \cmd{Operational
        Mode} (what it actually became).
  \item Either the VLAN does not exist in the VLAN database, or it exists and is
        shut down. \cmd{show vlan brief} distinguishes them.
\end{tightnum}

\subsection*{Chapter 9 --- Edge-Host Connectivity: PoE, Voice, APs and IoT}
\begin{tightnum}
  \item The voice VLAN, tagged. The phone learned its VLAN ID from CDP or LLDP-MED
        advertised by the switch. The PC's traffic stays untagged in the data
        VLAN, which is why one cable carries both.
  \item The switch has only 4.2\,W of power budget left, which is not enough for
        the AP. This is a \emph{budget} problem, not a port problem --- the port
        is willing (\cmd{auto}) and the supply cannot meet the request. Add a
        power supply, move the AP, or free capacity elsewhere.
  \item Because it tunnels all client traffic to the wireless controller inside
        CAPWAP. The SSID-to-VLAN mapping is done by the controller, so the switch
        port only needs to reach the controller. A \emph{standalone} AP bridges
        locally and does need a trunk.
\end{tightnum}

\subsection*{Chapter 10 --- Trunking, 802.1Q and SVIs}
\begin{tightnum}
  \item An access port. \cmd{dynamic auto} will accept a trunk if asked but never
        asks, so with neither side initiating, no trunk forms.
  \item No access port in that VLAN is up, so the SVI has nothing to be up for.
        Confirm with \cmd{show vlan brief} --- an SVI needs the VLAN to exist and
        at least one active port in it (or \cmd{no autostate}).
  \item Because a trunk that fails to come up is obvious and fails safe --- no
        traffic passes and somebody investigates. A native VLAN mismatch leaves
        the trunk up and quietly bridges two different VLANs together, merging
        broadcast domains that were meant to be separate. It is a security
        problem that looks like a working link.
\end{tightnum}

\subsection*{Chapter 11 --- EtherChannel and LACP}
\begin{tightnum}
  \item No channel forms --- both ends wait to be asked. The individual links stay
        up, so spanning tree sees parallel paths and blocks all but one. The
        result is a working but single-link connection, which is why the symptom
        is ``it works, but it is not four gigabits''.
  \item \cmd{(s)} is suspended: a parameter mismatch. Compare speed, duplex, and
        the allowed VLAN list and native VLAN on each member. \cmd{show
        interfaces <int> switchport} and \cmd{show interfaces status} show them.
  \item Because load balancing is per-\emph{flow}, not per-packet. A hash of
        addresses and ports selects one member for a conversation, and one file
        copy is one conversation. Bundling raises aggregate capacity across many
        flows, not the speed of any single one.
\end{tightnum}

\subsection*{Chapter 12 --- Rapid PVST+ and the Protection Features}
\begin{tightnum}
  \item The one ending \cmd{...0050}. With equal priorities the lowest MAC address
        wins the tie-break, and \cmd{0050} is numerically lowest.
  \item No, it is correct and necessary. An alternate port blocking is spanning
        tree doing exactly its job --- it has found a redundant path and is
        keeping it available without creating a loop.
  \item Because PortFast removes the delay that would otherwise let spanning tree
        detect a loop before the port forwards. If somebody plugs a switch into
        that port, a loop forms immediately. BPDU guard restores the safety by
        err-disabling the port the instant a BPDU appears --- which is the only
        thing that should never arrive on an access port.
\end{tightnum}

\subsection*{Chapter 13 --- CDP, LLDP and Documenting a Network}
\begin{tightnum}
  \item There is no switch between you and it that processes CDP --- so the
        intervening devices are hubs, media converters, or unmanaged switches
        that pass the frames through. Neighbour discovery only shows
        \emph{directly connected} devices, so a distant one appearing means the
        path is not what the diagram says.
  \item Because entries persist for the advertised holdtime, 180 seconds by
        default for CDP. Until it expires, the table describes a neighbour that
        is no longer there.
  \item \cmd{show cdp neighbors detail}. The summary form gives neither.
\end{tightnum}

\subsection*{Chapter 14 --- Troubleshooting Layer 2}
\begin{tightnum}
  \item Either the VLAN is missing from a trunk's allowed list, or it has been
        deleted or shut down on one switch. \cmd{show interfaces trunk} separates
        them --- read \cmd{Vlans allowed and active}, which reflects both the
        allowed list and whether the VLAN exists.
  \item To fix the cause. Err-disable is a symptom; whatever triggered it --- a
        port security violation, a BPDU on a PortFast port --- is still there and
        will trigger again the moment the port comes up.
  \item Because everything that still works is evidence you can trust, and it
        eliminates whole layers and paths at once. A user's report tells you one
        thing is broken; knowing that four similar things are fine tells you
        where the fault \emph{cannot} be, which is a much larger reduction in the
        search space.
\end{tightnum}

\section*{Part III --- IP Routing}
\addcontentsline{toc}{section}{Part III --- IP Routing}

\subsection*{Chapter 15 --- The Routing Table and How a Router Decides}
\begin{tightnum}
  \item The static \cmd{10.5.0.0/16}, because it is the longer prefix match.
        Administrative distance never enters into it --- longest prefix match is
        applied first and decides between \emph{entries}.
  \item The static route, distance 1 against OSPF's 110. The OSPF route is not
        installed in the routing table but remains in the OSPF database, so it
        will appear if the static route is removed or its next hop becomes
        unreachable.
  \item A layer 2 problem on the path; an ACL blocking the traffic; or no return
        route at the far end. \cmd{traceroute} distinguishes the most common by
        showing where the packets stop, and whether they get there and cannot get
        back.
\end{tightnum}

\subsection*{Chapter 16 --- Static Routing, IPv4 and IPv6}
\begin{tightnum}
  \item The next hop is not reachable. A router will not install a static route
        whose next-hop address it cannot resolve to an interface. Confirm with
        \cmd{show ip route <next-hop>} --- if that returns nothing, the static
        route has nowhere to point.
  \item Nothing will happen: 100 is \emph{lower} than OSPF's 110, so the ``backup''
        route wins permanently and the OSPF route is never installed. A floating
        static must have a distance \emph{higher} than the protocol it backs up.
  \item Because a link-local address is only unique within a link, and a router
        with several interfaces has no way to tell which one \cmd{fe80::2} refers
        to. Naming the exit interface removes the ambiguity.
\end{tightnum}

\subsection*{Chapter 17 --- OSPFv2, Single Area}
\begin{tightnum}
  \item When both routers are DROTHERs. Two non-designated routers on a broadcast
        segment deliberately stop at 2-WAY --- they exchange databases with the DR
        and BDR, not with each other. It is a fault only if one of them should
        have been the DR or BDR.
  \item \cmd{0.0.3.255}. The mask for \cmd{/22} is \cmd{255.255.252.0}, and
        $255-252 = 3$.
  \item The router ID was derived from that interface's address, so when the
        interface flaps the router ID changes and every adjacency rebuilds. Fix
        it with an explicit \cmd{router-id}, or a loopback, which never flaps.
\end{tightnum}

\subsection*{Chapter 18 --- OSPFv3 for IPv6}
\begin{tightnum}
  \item Because the router ID is a 32-bit value written in IPv4 format regardless
        of the address family, and it is normally derived from an IPv4 address on
        the device. With no IPv4 anywhere, none can be derived, and the process
        will not start at all --- the interface commands are accepted and do
        nothing.
  \item Nothing. They are separate processes with separate databases, neighbours
        and adjacencies. A working v2 adjacency proves the link and layer 2 are
        fine, which is useful, but says nothing about v3's configuration.
  \item Per interface, with \cmd{ipv6 ospf <pid> area <n>}. \emph{Every} IPv6
        prefix on that interface is advertised --- there are no network statements
        and no wildcard masks, which is why \cmd{passive-interface} matters more
        in OSPFv3 than in OSPFv2.
\end{tightnum}

\subsection*{Chapter 19 --- First Hop Redundancy: HSRP and VRRP}
\begin{tightnum}
  \item Because the virtual MAC address moves with the virtual IP. The host's ARP
        cache still maps the gateway IP to the same MAC; the new active router
        simply starts answering to it. The host never learns anything happened,
        which is the entire point.
  \item That the two routers are not seeing each other's hellos --- so each
        believes it is alone. Causes: mismatched group numbers, different
        subnets, a VLAN or trunk fault between them, or an ACL blocking the
        multicast. Both then claim active, and two devices answer for the same IP.
  \item The priority drops to 140, which is still above the peer's 100, so
        \textbf{nothing happens} --- the router keeps the active role despite
        having lost the tracked interface. That is almost certainly not what was
        intended: the decrement must be large enough to take the priority
        \emph{below} the peer's, so 51 or more.
\end{tightnum}

\section*{Part IV --- Network Services and Security}
\addcontentsline{toc}{section}{Part IV --- Network Services and Security}

\subsection*{Chapter 20 --- DHCP: Client, Server and Relay}
\begin{tightnum}
  \item Discover, Offer and Request are broadcast; only the Ack is unicast. If the
        Request were unicast, a second server that had also made an offer would
        never learn that its offer was declined, and would hold that address
        reserved for nothing until it timed out.
  \item It establishes only that the host asked and \textbf{nothing answered}. It
        does not establish where the failure is --- it could be the switch port,
        the VLAN, the SVI, the relay, or the server, and it says nothing about
        whether the server is healthy.
  \item On the interface facing the \textbf{clients} --- the SVI in the subnet
        where the broadcasts are. The relay writes its own address on that
        interface into \cmd{giaddr}, which tells the server which pool to use and
        gives it a unicast address to reply to.
  \item A relay or broadcast-path fault. Existing hosts renew at T1 by
        \emph{unicast} directly to the server, which needs no relay, so they
        survive. Only newly booting hosts need the broadcast path, so only they
        fail.
  \item No. DHCP delivered a correct address and a correct gateway, which is its
        entire job. The fault is elsewhere --- routing, NAT, an ACL, or DNS. Check
        the DNS servers the lease carried before looking further afield.
\end{tightnum}

\subsection*{Chapter 21 --- DNS, and Diagnosing Records}
\begin{tightnum}
  \item That the network path and the server are fine and the fault is name
        resolution. Next, query the record explicitly and against a known server:
        \cmd{nslookup -type=A www.example.edu 8.8.8.8}, and compare with what your
        own resolver returns.
  \item \cmd{NXDOMAIN} means the \emph{name} does not exist. \cmd{NOERROR} with an
        empty answer means the name exists but has no record \emph{of the type you
        asked for}. They look identical to a user and mean entirely different
        things.
  \item Because a CNAME means ``this name is an alias for that name'', and a name
        that is an alias cannot also carry data of its own. Standards prohibit
        other records alongside it, and mail servers differ in how they cope ---
        which produces the worst symptom, where mail from some senders arrives
        and from others does not.
  \item AAAA. A dual-stack client prefers IPv6 and will try the AAAA address
        first; if it is wrong or unreachable the client waits for a timeout before
        falling back to IPv4. IPv4-only clients never see it, so the site is slow
        for some and instant for others.
  \item Usually whoever owns the address block --- your ISP or cloud provider, not
        you, because reverse DNS is delegated separately in
        \cmd{in-addr.arpa}. Without it, outbound mail is commonly rejected or
        scored as spam, because receiving servers check that the connecting
        address reverse-resolves and forward-resolves back.
\end{tightnum}

\subsection*{Chapter 22 --- NAT and PAT on IOS XE}
\begin{tightnum}
  \item Inside local \cmd{10.1.1.9}; inside global \cmd{198.51.100.4}; outside
        global \cmd{203.0.113.80}; outside local \cmd{203.0.113.80} --- the same,
        because plain NAT does not translate the outside host's address.
  \item \cmd{overload}. It translates the port as well as the address, so one
        public address serves many hosts. Without it, dynamic NAT lends addresses
        one-to-one and \textbf{silently fails} when the pool is exhausted.
  \item It rules out a pool or ACL problem --- misses would be non-zero if packets
        had been considered and refused. Both counters at zero means no packet was
        ever evaluated for translation, which points at the missing \cmd{ip nat
        inside} or \cmd{ip nat outside} interface markers.
  \item The \textbf{inside global} address, the translated one. An inbound ACL on
        the outside interface is evaluated \emph{before} the packet is
        un-translated, so it sees the address that is on the wire at that
        interface.
  \item Because it is an address-conservation mechanism. It happens to make
        unsolicited inbound connections difficult, since no translation entry
        matches them, but that is a side effect rather than a policy --- it cannot
        be audited, it states no intent, and it is trivially bypassed by anything
        the inside host initiates. Use an ACL or firewall for policy.
\end{tightnum}

\subsection*{Chapter 23 --- Device Access, AAA and Secure File Transfer}
\begin{tightnum}
  \item \cmd{local} is used only when the TACACS+ group \emph{does not respond at
        all}. If a server answers and rejects the credentials, that is an answer,
        the list ends, and the local database is never consulted. Fallback
        protects against a dead server, not a forgotten password.
  \item TACACS+ is TCP 49 and encrypts the entire packet body; RADIUS is UDP
        1812/1813 and encrypts only the password field. Choose \textbf{TACACS+}
        for switch administrators --- it separates authorization fully and
        supports per-command control.
  \item Because the RSA key pair is named \cmd{hostname.domain}. Without a domain
        name the device has no full name to label the key with and refuses.
  \item \cmd{aaa new-model}, an \cmd{aaa authentication login} method, and
        \cmd{aaa authorization exec} --- plus \cmd{ip scp server enable} itself.
        Without the authorization line, users authenticate successfully and every
        transfer is then refused.
  \item TFTP has no user authentication, sends everything in clear text, and runs
        over UDP with no integrity protection. A configuration file in flight ---
        including its shared secrets --- is readable by anyone on the path, and
        anyone who can reach the server can fetch any file whose name they guess.
\end{tightnum}

\subsection*{Chapter 24 --- Access Control Lists}
\begin{tightnum}
  \item Mask \cmd{255.255.240.0}, so the wildcard is \cmd{0.0.15.255}. It matches
        \cmd{10.20.16.0} through \cmd{10.20.31.255} --- sixteen consecutive /24s.
  \item Because a standard ACL can only match the \emph{source}. Placed near the
        source it would block that host from everything, not just the one
        destination. An extended ACL can name the destination too, so it is
        precise enough to sit near the source, which saves carrying doomed traffic
        across the network.
  \item That everything is being dropped by an earlier line --- no packet from
        anywhere reached the bottom of the list. An ACL whose final permit never
        matches is denying all traffic it sees.
  \item TCP segments with the ACK or RST flag set --- replies within a conversation
        somebody inside already started. It is not stateful because it inspects
        \emph{flags} on each packet independently and keeps no record of
        connections; a crafted packet with ACK set will match.
  \item \textbf{DHCP} --- without it hosts get no address and the VLAN looks
        completely dead. \textbf{DNS} --- everything works by IP and nothing by
        name. \textbf{ICMP} --- users and you lose the ability to test anything,
        and troubleshooting becomes guesswork.
\end{tightnum}

\subsection*{Chapter 25 --- Layer 2 Security}
\begin{tightnum}
  \item Port security --- one device claiming many MAC addresses or an
        unauthorised device. DHCP snooping --- a rogue DHCP server. Dynamic ARP
        inspection --- ARP poisoning. Storm control --- flooding. RA guard --- a
        rogue IPv6 router.
  \item Because DAI validates ARP against the DHCP snooping \textbf{binding
        table}, which records which MAC legitimately holds which IP on which port.
        Without snooping running there is no table, and DAI has nothing to
        validate against.
  \item It drops offending traffic silently --- no counter, no log, no trap. A
        user whose traffic is being discarded reports ``the network is slow'' and
        nothing anywhere records why. Prefer \cmd{restrict}, which drops and
        reports.
  \item Because the phone appears in the voice VLAN, the PC in the data VLAN, and
        the phone commonly registers a second address of its own. A maximum of 2
        produces violations that come and go with phone firmware.
  \item It is statically addressed, so it never took a DHCP lease and has no entry
        in the snooping binding table. DAI therefore drops every ARP reply it
        sends, so nothing can resolve its MAC address and it becomes unreachable
        while remaining healthy. Fix with an ARP ACL naming its address and MAC,
        applied with \cmd{ip arp inspection filter}.
  \item The uplink toward the distribution layer, for \textbf{DHCP snooping} and
        for \textbf{dynamic ARP inspection}. Untrusted uplinks take out the whole
        switch in both cases.
\end{tightnum}

\subsection*{Chapter 26 --- IPsec VPNs}
\begin{tightnum}
  \item ESP, IP protocol \textbf{50}. AH provides integrity and authentication but
        no encryption.
  \item Tunnel mode encrypts the entire original packet and adds a new IP header;
        transport mode keeps the original header and encrypts only the payload. A
        site-to-site VPN uses \textbf{tunnel mode}.
  \item Because AH authenticates the entire packet including the source and
        destination addresses in the IP header, and NAT's whole function is to
        rewrite them. The integrity check therefore fails by definition and the
        far end discards the packet. No configuration fixes this.
  \item Phase 1 authenticates the peers and builds a secure channel between them
        --- the IKE security association. Phase 2 uses that channel to negotiate
        the security associations that protect user traffic, including which
        traffic is protected.
  \item Because plain IPsec tunnel mode carries unicast IP only, and OSPF depends
        on multicast. GRE can carry multicast, so the routing protocol runs inside
        a GRE tunnel, which IPsec then protects.
  \item Path MTU. The added headers push large packets past the MTU somewhere on
        the path, so small packets pass and large ones are dropped or fragmented
        badly. A good reminder that ``it pings'' is not ``it works''.
\end{tightnum}

\section*{Part V --- Wireless, Virtualisation and Operations}
\addcontentsline{toc}{section}{Part V --- Wireless, Virtualisation and Operations}

\subsection*{Chapter 27 --- Wireless Principles}
\begin{tightnum}
  \item Because channels are numbered 5\,MHz apart but each is 20\,MHz wide, so
        five channel numbers of separation are needed to avoid overlap. In a band
        of eleven usable channels that gives exactly three: \textbf{1, 6 and 11}.
  \item \textbf{Adjacent-channel} is worse. Two APs on the same channel can hear
        each other, so CSMA/CA works and they share the airtime --- throughput
        halves but nothing is corrupted. On partly overlapping channels neither
        can decode the other, so neither defers, both transmit, and frames are
        destroyed and retransmitted.
  \item The \textbf{noise floor}, so you can calculate SNR. A signal of $-62$ is
        strong; if the noise floor is around $-70$ the SNR is only 8\,dB, well
        below the 25\,dB voice needs. That would make the report entirely normal.
  \item Personal shares one passphrase among everyone and identifies nobody;
        Enterprise authenticates each user individually via 802.1X against a
        RADIUS server, so one account can be revoked without touching anyone else.
        WPA2 uses AES-CCMP with a 4-way handshake that can be captured and
        attacked offline; WPA3 replaces it with SAE, which resists offline attack
        and adds forward secrecy.
  \item Bodies are largely water, which absorbs 2.4\,GHz strongly --- so a full
        room attenuates far more than an empty one. And a full room means many
        clients contending for the same airtime on a half-duplex shared medium.
  \item Non-Wi-Fi interference --- a microwave oven, a video sender, a cordless
        phone. It transmits no frames, so it never appears in a Wi-Fi scan; it
        simply raises the noise floor. Finding it needs a \textbf{spectrum
        analyser}, which looks at raw RF energy rather than decoded frames.
\end{tightnum}

\subsection*{Chapter 28 --- Troubleshooting Client Connectivity}
\begin{tightnum}
  \item (1) Do I have an address? Failure points to DHCP or layer 2. (2) Ping my
        gateway --- failure points to the local segment. (3) Ping something remote
        by IP --- failure points to routing, NAT or an ACL. (4) Ping the same thing
        by name --- failure points to DNS and nothing else.
  \item \cmd{ipconfig /all} on Windows, \cmd{ifconfig} on macOS, \cmd{ip addr} on
        Linux. Note that \cmd{ifconfig} is deprecated on Linux and normal on
        macOS.
  \item It can reach hosts on its own subnet and nothing else. Everything local
        works, so the user may report it as ``the internet is down'' rather than
        as a network fault.
  \item Because the mask decides which destinations the host believes are local.
        With a mask that is too short, the host ARPs directly for addresses that
        are actually on other subnets, gets no answer, and fails --- while
        destinations genuinely outside the over-wide range still use the gateway
        correctly and work.
  \item Either it is being placed in a VLAN with no DHCP scope --- often because
        MAC randomisation means it is not recognised by a MAC-based authorisation
        list --- or the relay or scope for its VLAN is broken. Tell them apart by
        checking whether \emph{other} clients on the same SSID and VLAN are
        leasing normally; if they are, it is specific to this device.
  \item The locally-administered bit is set --- the second-lowest bit of the first
        octet. In practice the first octet's second hex digit will be 2, 6, a or e
        (as in \cmd{8A}), and the address will not match the device's label or
        the manufacturer OUI.
\end{tightnum}

\subsection*{Chapter 29 --- Hypervisors, Virtual Machines and Containers}
\begin{tightnum}
  \item One layer. A type~1 hypervisor runs directly on the hardware and \emph{is}
        the operating system; a type~2 runs as an application on top of a host
        operating system.
  \item The \textbf{kernel}. Consequences: containers are far smaller and start in
        milliseconds because there is no operating system to boot; and isolation
        is weaker, because a kernel vulnerability is shared by every container on
        the host.
  \item Because the VMs on it live in several VLANs, and their frames arrive at
        the physical switch tagged. An access port would carry one VLAN and
        discard the rest, stranding every VM outside it.
  \item Because a hypervisor presents many MAC addresses on one port --- one per
        vNIC plus the host's own --- and they change as machines start, stop and
        migrate. Port security will exceed its maximum and err-disable the port,
        taking down every VM on the host.
  \item Traffic between two VMs on the same host and VLAN. It is switched inside
        the hypervisor's virtual switch and never leaves the server, so it never
        reaches the physical switch to be seen, filtered or captured there.
  \item Live migration. A running VM was moved to another host, so its MAC address
        legitimately appears on a different switch port seconds later. It looks
        exactly like a loop and is not.
\end{tightnum}

\subsection*{Chapter 30 --- Management Approaches, SNMP and Syslog}
\begin{tightnum}
  \item Device-based, controller-based, cloud-based, automation-based, and
        infrastructure as code. The test: \textbf{does the network still have a
        management system when the tool is switched off?} For controller-based,
        yes --- and you have just lost it. For automation-based, no --- nothing
        about the devices changed.
  \item Facility \cmd{\%PM}, severity \textbf{4}, mnemonic \cmd{ERR\_DISABLE}.
        Severity 4 is \emph{warning}.
  \item \cmd{warnings} is severity 4, so it records levels 0 to 4 and
        \textbf{discards} 5, 6 and 7 --- notification, informational and
        debugging.
  \item A trap is sent once over UDP and never acknowledged, so it can be lost
        silently. An inform is acknowledged and retransmitted until it is.
        Use informs for anything you would be sorry to miss, traps for
        high-volume low-consequence events.
  \item Because the community string travels in \textbf{clear text} in every
        polling packet, hundreds of times an hour. Anyone who can capture one
        packet has it, so obscurity is irrelevant. A read-write string is
        equivalent to an enable password.
  \item It sets the \emph{severity threshold for messages sent to a syslog
        server}. It has nothing whatever to do with SNMP traps, despite the name.
\end{tightnum}

\subsection*{Chapter 31 --- Ansible and Configuration Management}
\begin{tightnum}
  \item It installs nothing on the managed device and connects over SSH. It
        decides the choice because you cannot install an agent on a Catalyst
        switch --- so any agent-based tool (Puppet, Chef) is simply unavailable
        for network hardware, while Ansible needs only the SSH access already
        configured.
  \item An operation produces the same result whether applied once or many times.
        In the recap you see it as \cmd{ok} rather than \cmd{changed} on a device
        that is already in the desired state --- so a second run of the same
        playbook reports \cmd{changed=0}.
  \item \cmd{ios\_command} runs EXEC commands and returns output; it cannot make a
        configuration change. \cmd{ios\_config} enters configuration mode and
        applies lines. A playbook reporting success while changing nothing is very
        often the former where the latter was meant.
  \item It writes the running configuration to startup-config, but only when that
        run actually changed something. Without it the change lives in the running
        configuration only and is lost at the next reload --- which shows up weeks
        later as a fault correlating with reboots rather than with the change.
  \item \cmd{unreachable} means Ansible could not log in to the device at all.
        \cmd{failed} means it logged in and the task did not work. They point at
        completely different problems.
  \item \cmd{--check} is a dry run --- it reports what it would do and touches
        nothing. \cmd{--diff} shows the exact lines. Omit them only when you have
        already run with them and are applying the change you just reviewed.
\end{tightnum}

\subsection*{Chapter 32 --- AI in Network Operations}
\begin{tightnum}
  \item Generative AI produces content in response to a prompt and stops; agentic
        AI pursues a goal, uses tools to act, and evaluates its own results.
  \item Data classification, output format, persona and instructions.
        \textbf{Data classification} is listed first, because a prompt that
        discloses credentials or real addressing is wrong regardless of how good
        the rest of it is --- and the disclosure cannot be undone.
  \item Credentials of any kind (passwords, hashes, community strings, pre-shared
        keys, tokens); personal data; and real public IP addressing or anything
        the organisation classifies as confidential.
  \item How much breaks if the action is wrong. It governs autonomy because the
        question is not how capable the system is but how bad a confident mistake
        would be --- clearing one err-disabled port is reversible and local, while
        changing a core route map reaches the whole organisation.
  \item A hallucination. The defence is verification: check the command exists on
        your platform, check it against the exhibit or the device, and keep change
        control. Fluency is not evidence.
  \item Any three of: anomaly detection, event correlation and noise reduction,
        root-cause suggestion, configuration drafting and review, documentation
        and summarisation, capacity forecasting, guided remediation.
\end{tightnum}

\section*{Part VI --- Sitting the Exam}
\addcontentsline{toc}{section}{Part VI --- Sitting the Exam}

\subsection*{Chapter 33 --- A Troubleshooting Method for the 28 Per Cent}
\begin{tightnum}
  \item Define, gather, analyse, eliminate, hypothesise, test, solve. Steps
        \textbf{5 and 6} form the loop --- a hypothesis whose prediction fails is
        discarded and replaced, not patched.
  \item Bottom-up for a suspected physical or hardware fault, or when nothing
        works. Top-down when one application fails and everything else is fine.
        Divide and conquer as the default when you have no strong prior. Follow
        the path when two specific points cannot reach each other. Compare working
        to broken whenever an identical working case exists.
  \item \emph{One user or many?} --- sets the whole search space. \emph{Has it ever
        worked?} --- separates a build fault from a change fault. \emph{What
        changed?} --- finds the cause of most faults in previously working
        networks. \emph{Does it fail by name but work by address?} --- answers
        ``is this DNS'' in one test.
  \item Because it converts a check into information whichever way it goes. Without
        a prediction, a fix that does not resolve the symptom reads as ``the fix
        did not work'' rather than ``the hypothesis was wrong'', and you re-check
        something already correct.
  \item Either the hypothesis was wrong, or there is more than one fault. Go back
        to scoping rather than re-checking the thing you just verified --- and
        keep the fix if it corrected a genuine defect, but record it separately.
  \item The precise symptom and who is affected; what has been ruled out and the
        evidence for it; what has been changed, including anything not yet undone;
        and the current hypothesis with the next test.
\end{tightnum}

\subsection*{Chapter 34 --- Exam Strategy}
\begin{tightnum}
  \item \textbf{You cannot return to a previous question.} It invalidates
        skimming the paper first, leaving hard questions to come back to, and
        flagging anything for review.
  \item 60 to 75 seconds. At 90 seconds, eliminate what you can, choose the best
        remaining option, and move on --- the question you have not seen yet may be
        one you would answer in twenty seconds.
  \item Because you cannot come back to it, so a deferred simulation is an
        abandoned one. Read the whole task first; plan before typing; configure
        only what is asked; verify with a \cmd{show} command --- and save if the
        task says to.
  \item The \textbf{question} first. Otherwise you study an entire configuration
        or topology to answer something about one line of it.
  \item The \emph{describe} topics --- wireless principles, virtualisation, VPNs,
        management approaches, AI --- because they are recall-heavy and recall
        decays fastest. Configuration skill does not.
  \item Cisco publishes the 120-minute duration, the five domains and their
        weights, and the 29 topics with their wording. It does \textbf{not}
        publish the number of questions or the passing score.
\end{tightnum}
